\documentclass[11pt]{amsart}
\usepackage[T1]{fontenc}
\usepackage{lmodern,microtype,amsmath,amssymb,amsthm,mathtools}
\usepackage[margin=1.1in]{geometry}
\usepackage[hidelinks]{hyperref}
\usepackage{xurl}
\usepackage{booktabs}

\newtheorem{theorem}{Theorem}
\newtheorem{corollary}[theorem]{Corollary}
\newtheorem{lemma}[theorem]{Lemma}
\newtheorem{proposition}[theorem]{Proposition}
\theoremstyle{remark}
\newtheorem{remark}[theorem]{Remark}
\numberwithin{equation}{section}

\newcommand{\T}{\mathbb{T}}
\newcommand{\PP}{\mathbb{P}}
\newcommand{\E}{\mathbb{E}}
\newcommand{\dd}{\mathrm{d}}

\hypersetup{pdftitle={Frozen Percolation with Rare Lightning up to the Critical Time},pdfauthor={Yitzchak Shmalo}}

\begin{document}

\title[Rare lightning up to the critical time]{Frozen Percolation with Rare Lightning up to the Critical Time}
\author{Yitzchak Shmalo}
\address{Einstein Institute of Mathematics, The Hebrew University of Jerusalem, Jerusalem, Israel}
\email{yitzchak.shmalo@gmail.com}
\date{}
\subjclass[2020]{60K35, 82B43, 60J80}
\keywords{frozen percolation, lightning, forest fire, three-regular tree, recursive distributional equation, local limit}

\begin{abstract}
In frozen percolation on a finite graph, edges open at independent
uniform times, vertices are struck by lightning at rate $\lambda$, and a
strike freezes the open cluster of the struck vertex. On graphs of
maximum degree at most three, the probability that a given vertex is
frozen by time $1/2$, the critical time of the three-regular tree, is at
most
$\frac34\lambda\log(1/\lambda)+O(\lambda)$, uniformly in the graph. So
along such graphs converging locally to the three-regular tree with
$\lambda_n\to0$, the process on $[0,1/2]$ converges to Aldous's frozen
percolation on the tree, where clusters freeze on becoming infinite,
none before time $1/2$. On the tree itself, an exact recursion for the
freezing time shows that the root is frozen by time $1/2$ with
probability
$\frac34\lambda[\log(1/\lambda)-\log\log(1/\lambda)+\log4-\frac13+o(1)]$,
so the constant $\frac34$ in the bound cannot be lowered. Convergence after time $1/2$ remains
open.
\end{abstract}

\maketitle

\section{Introduction}

Frozen percolation was introduced by Aldous~\cite{Ald00}. Edges of the
three-regular tree $\T$ open at independent uniform times in $[0,1]$, and
an open cluster freezes at the moment it becomes infinite, after which no
edge touching it opens. Aldous constructed such a process on $\T$; we
call it Aldous's process. Ráth, Swart and Terpai~\cite{RST21} showed
that whether an edge of $\T$ freezes in Aldous's process is not a
function of the edge times, and they asked how the process arises from
finite graphs. Infinite clusters do not exist on a finite graph, so they freeze
clusters by lightning instead.

\smallskip
\noindent\textbf{The model.} Let $G=(V,E)$ be a finite graph and let
$\lambda>0$. Let $(U_e)_{e\in E}$ be independent and uniform on $[0,1]$,
and let the \emph{lightning times} $(\Lambda_v)_{v\in V}$ be independent exponential random variables
with mean $\lambda^{-1}$, independent of the edge times. Every vertex is
\emph{warm} at time $0$. At time $U_e$ the edge $e$ opens, provided that
neither endvertex is frozen. At time $\Lambda_v$ every vertex of the open
cluster of $v$ becomes frozen. Frozen vertices stay frozen and open edges
stay open. We write $\tau_v$ for the time at which $v$ becomes frozen
($\tau_v=\infty$ if it never does). This is the finite-graph model of
\cite[Section~1.7]{RST21}, where $\lambda$ is called the lightning rate.
Replacing the single time $\Lambda_v$ by the points of a Poisson process
of rate $\lambda$ changes nothing, since strikes at frozen vertices have
no effect. We call the edge times and the lightning times the \emph{marks}.

Ráth, Swart and Terpai consider graphs $G_n$ with $n$ vertices that
converge locally to a limit graph, with rates satisfying
$n^{-1}=o(\lambda_n)$ and $\lambda_n\to0$. They pose the following
problem~\cite[Problem~2]{RST21}: show that frozen percolation on the
three-regular tree is the local weak limit of frozen percolation on a
suitable sequence of finite graphs. They also observe
that up to the critical time $p_c$ of the limit graph, open clusters
grow as in ordinary percolation \cite[Section~1.7]{RST21}; for $\T$ the
critical time is $1/2$. Let $C_s(v)$ be the cluster of $v$ in the graph
of edges with $U_e\le s$. A vertex frozen at time $s$ has a struck vertex
in $C_s(v)$, so $\PP(\tau_v\le t)\le\lambda\int_0^t\E|C_s(v)|\,\dd s$,
and on graphs of maximum degree at most three
$\E|C_s(v)|\le(1+s)/(1-2s)$. This bound is finite for $t<1/2$ and
infinite at $t=1/2$: on $\T$ the expected cluster size equals
$(1+s)/(1-2s)$, which is not integrable at $1/2$. Near the critical time
we replace the expectation by the generating function $\psi$ of
the total progeny of a critical Galton--Watson tree whose root has
$\mathrm{Bin}(3,1/2)$ children and whose other vertices have
$\mathrm{Bin}(2,1/2)$ children, given in
\eqref{eq:A1} below. An upper bound of this kind does not show that its
constant is attained. On $\T$, removing an edge at the root leaves a
rooted binary tree, a \emph{branch}; the subtrees below a vertex are
independent branches, and the freezing time of the root of a branch
satisfies an exact recursive distributional equation of the kind
surveyed by Aldous and Bandyopadhyay~\cite{AB05}. Written for the
edges, it is the recursion of Aldous's process as given by Ráth, Swart
and Terpai~\cite[Section~1.2]{RST21}, with the lightning time of the root
of the branch added to the minimum. It reduces to a
differential equation, and the logarithm that this equation produces at
$t=1/2$ shows that the constant is attained.

We make the observation of Ráth, Swart and Terpai quantitative and compute the first-order
effect of the lightning before time $1/2$ on the tree. The first result
is a bound that depends on the graph only through its maximum degree.

\begin{theorem}\label{thm:A}
Let $G$ be a finite graph of maximum degree at most three, let $v$ be a
vertex of $G$, and let $\lambda>0$. For every $0<\varepsilon<1/2$,
\begin{equation}\label{eq:A1}
\PP(\tau_v\le\tfrac12)\le
\lambda\Bigl[-\frac14+\frac\varepsilon2+\frac34\log\frac1{2\varepsilon}\Bigr]
+1-\psi\bigl(e^{-\lambda\varepsilon}\bigr),
\qquad
\psi(z)=\frac{1-\sqrt{1-z}}{\bigl(1+\sqrt{1-z}\bigr)^2}.
\end{equation}
In particular, for $0<\lambda\le1$,
\begin{equation}\label{eq:A2}
\PP(\tau_v\le\tfrac12)\le \beta(\lambda):=
\frac{3\lambda}{4}\log\frac2\lambda+\frac{5\lambda}{4}+\frac{\lambda^2}{8}.
\end{equation}
\end{theorem}

The bound needs no assumption on $n\lambda$ and holds for every vertex,
not only for a typical one. It gives the convergence of the process up to
time $1/2$. For a finite graph $G$ and a vertex $o$ let $B_r(G,o)$ be the
ball of radius $r$ around $o$, and let $\T_r$ be the ball of radius $r$
in $\T$, which has $3\cdot2^r-2$ vertices. Up to time $1/2$, Aldous's
process on $\T$ freezes nothing, since critical percolation on $\T$ has
no infinite cluster; every edge $e$ is open from time $U_e$ on.

\begin{corollary}\label{cor:B}
Let $G_n$ be finite graphs of maximum degree at most three, let $o_n$ be
a uniformly chosen vertex of $G_n$, and for $r\ge1$ let
$b_n(r)=\PP(B_r(G_n,o_n)\not\cong\T_r)$. Let $\lambda_n\in(0,1]$. For
every $r\ge1$ there is a coupling of the frozen percolation process on $G_n$ with rate
$\lambda_n$, rooted at $o_n$, with Aldous's process on $\T$, under which,
outside an event of probability at most
\begin{equation}\label{eq:B1}
b_n(r)+(3\cdot2^r-2)\,\beta(\lambda_n),
\end{equation}
the two balls of radius $r$ are identified and every edge and every
vertex of the ball has the same open/closed and warm/frozen history on
$[0,1/2]$ in both processes. In particular, if $G_n$ converges locally to
$\T$, that is $b_n(r)\to0$ for every $r$, and $\lambda_n\to0$, then for
every $r$ the history on $[0,1/2]$ of the ball of radius $r$ around $o_n$
converges in law to the history of $\T_r$ in Aldous's process. On
$\T$ that restriction has no frozen vertex almost surely, so on
$[0,1/2]$ the limit is ordinary percolation.
\end{corollary}

Theorem~\ref{thm:A} and Corollary~\ref{cor:B} are proved in
Section~\ref{sec:A}. The main result shows that the order
$\lambda\log(1/\lambda)$ and the constant $3/4$ are correct. On $\T$ the
process up to time $1/2$ is defined by running it on each critical
percolation cluster, all of which are finite.

\begin{theorem}\label{thm:C}
Let $o$ be the root of the three-regular tree $\T$ and $L=\log(1/\lambda)$.
For the frozen percolation process with lightning rate $\lambda$ on $\T$,
as $\lambda\downarrow0$,
\begin{equation}\label{eq:C1}
\PP(\tau_o\le\tfrac12)=\frac{3\lambda}{4}\Bigl[L-\log L+\log4-\frac13+o(1)\Bigr].
\end{equation}
\end{theorem}

Balls $\T_R$ rooted at their centre are finite graphs of maximum degree
three, and their root freezing probabilities tend to that of $\T$ as
$R\to\infty$ (Section~\ref{sec:sharp}). Together with
Theorem~\ref{thm:A} this gives
\begin{equation}\label{eq:sharp}
\lim_{\lambda\downarrow0}\ \sup_{G,v}\
\frac{\PP_G(\tau_v\le\frac12)}{\lambda\log(1/\lambda)}=\frac34,
\end{equation}
the supremum running over finite graphs of maximum degree at most three
and their vertices. The centred balls are not themselves a sequence
converging locally to $\T$, since about half of their vertices are
leaves; \eqref{eq:sharp} concerns the uniform bound only.

The proof of Theorem~\ref{thm:C} rests on the exact recursion for the
freezing time $X$ of the root of a branch (Section~\ref{sec:rec}). It
reduces to the ordinary differential equation
\[
h'=\lambda\,\frac{(1+h)(1+h-t)}{1-2t+h},\qquad h(0)=0,
\]
on $[0,1/2]$, where $h$ is an explicit function of the law of $X$
(Lemma~\ref{lem:ode}). As $\lambda\to0$ the denominator nearly
vanishes at $t=1/2$, and the logarithmic singularity of the integral of
$1/(1-2t+h)$ produces the term $\log(1/\lambda)$ (Section~\ref{sec:asym}).

\smallskip
\noindent\textbf{Relation to other work.}
Aldous constructed his process in~\cite{Ald00}, recursive
distributional equations of this kind are surveyed by Aldous and
Bandyopadhyay~\cite{AB05}, and Ráth, Swart and Terpai proved that the
process is not endogenous~\cite{RST21}. Van den Berg, Kiss and
Nolin~\cite{BKN12} study a percolation process on the planted binary
tree, the binary tree whose root has a single child, in which clusters freeze as soon as their size reaches a parameter
$N$, instead of by lightning, and show that as $N\to\infty$, at each time $t$, the probability that
the open cluster of the root equals a given finite cluster converges to
the corresponding probability in Aldous's process, and so does the
probability that this cluster is frozen~\cite[Theorem~1]{BKN12}. Corollary~\ref{cor:B} differs in three ways: it
concerns finite graphs converging to $\T$ rather than the tree itself,
clusters freeze by lightning rather than at a size threshold, and it
covers only times up to $1/2$. On the complete graph, frozen
percolation was introduced by Ráth~\cite{Rat09}, following a suggestion
of Aldous~\cite[Section~5.5]{Ald00}, and the related forest fire model by
Ráth and Tóth~\cite{RT09}. Ráth, Szőke and Yeo~\cite{RSY26} study the
complete graph with lightning rate $\lambda_n$, where struck components
are deleted. Their deterministic window curve is $w(t)=tq(t)$, where
$q$, a deterministic approximation of the fraction of vertices not yet
deleted, solves $q'=-\lambda q/(1-tq)$ with $q(0)=1$. At the critical
time $t=1$ the deleted fraction $1-q(1)=1-w(1)$ in this approximation is
of order $\lambda\log(1/\lambda)$, as they note
without proof \cite[Remark~2.18]{RSY26}; their Claim~2.17 gives the time
at which the window curve reaches criticality. Theorem~\ref{thm:C} is a tree
counterpart of this, with the constant $3/4$ identified.

Problem~2 of \cite{RST21} remains open. After time $1/2$ the
percolation clusters of $\T$ can be infinite, the construction by finite
clusters fails, and freezing can be transmitted into a ball from far
away.
Section~\ref{sec:sharp} proves \eqref{eq:sharp} and gives two examples,
and Section~\ref{sec:num} compares \eqref{eq:C1} with a numerical
solution of the differential equation.

\section{A uniform bound before the critical time}\label{sec:A}

Couple the process with ordinary percolation $E_t=\{e:U_e\le t\}$ by
using the same edge times. Open edges at time $t$ belong to $E_t$. Let
$C_t(v)$ be the cluster of $v$ in $(V,E_t)$, as in the introduction. A vertex freezes only at a
lightning time, and only if the struck vertex lies in its open cluster.
Hence
\begin{equation}\label{eq:necessary}
\tau_v=s\quad\Longrightarrow\quad \Lambda_u=s\ \text{ for some } u\in C_s(v).
\end{equation}

\begin{lemma}\label{lem:cluster}
Let $G$ be a finite or countably infinite graph of maximum degree at
most three. For $0\le t<1/2$,
\[
\E|C_t(v)|\le1+\frac{3t}{1-2t}=-\frac12+\frac{3}{2(1-2t)}.
\]
At $t=1/2$, $|C_{1/2}(v)|$ is stochastically dominated by the total
progeny $N$ of a Galton--Watson tree whose root has a
$\mathrm{Bin}(3,1/2)$ number of children and whose other vertices have
$\mathrm{Bin}(2,1/2)$ children, and $\E z^N=\psi(z)$ for
$0\le z\le1$.
\end{lemma}

\begin{proof}
Explore $C_t(v)$ breadth first, testing each unexplored edge at a
discovered vertex once. The root has at most three edges and every other
discovered vertex at most two unexplored edges, each open independently
with probability $t$; an edge leading to an already discovered vertex adds
nothing. So $|C_t(v)|$ is dominated by the total progeny of a
Galton--Watson tree with root offspring $\mathrm{Bin}(3,t)$ and other
offspring $\mathrm{Bin}(2,t)$, whose mean is $1+3t\sum_{k\ge0}(2t)^k$ for
$t<1/2$. For $t=1/2$, the progeny generating function $f$ of a
$\mathrm{Bin}(2,1/2)$ vertex satisfies $f=z\bigl((1+f)/2\bigr)^2$. With
$y=\sqrt{1-z}$ the root of this equation in $[0,1]$ is
$f=(1-y)/(1+y)$, and the root has
$\E z^N=z\bigl((1+f)/2\bigr)^3=z/(1+y)^3=(1-y)/(1+y)^2$.
\end{proof}

\begin{proof}[Proof of Theorem~\ref{thm:A}]
Put $t_0=1/2-\varepsilon$. By \eqref{eq:necessary}, $\tau_v\le t_0$
requires some $u$ with $\Lambda_u\le t_0$ and $u\in C_{\Lambda_u}(v)$. The
lightning times are independent of the edge times and have density at
most $\lambda$, so
\[
\PP(\tau_v\le t_0)\le\sum_u\int_0^{t_0}\lambda\,\PP(u\in C_s(v))\,\dd s
=\lambda\int_0^{t_0}\E|C_s(v)|\,\dd s
\le\lambda\Bigl[-\frac{t_0}2+\frac34\log\frac1{1-2t_0}\Bigr]
\]
by Lemma~\ref{lem:cluster}; this is the first term of \eqref{eq:A1}.

If $t_0<\tau_v\le1/2$, then by \eqref{eq:necessary} some vertex of
$C_{1/2}(v)\supseteq C_{\tau_v}(v)$ has its lightning time in
$(t_0,1/2]$. Given $|C_{1/2}(v)|=k$, the lightning times of these $k$
vertices are independent of the cluster, and the probability that one of
them falls in $(t_0,1/2]$ is
\[
1-\bigl(1-e^{-\lambda t_0}(1-e^{-\lambda\varepsilon})\bigr)^k
\le1-e^{-\lambda\varepsilon k}.
\]
The right side increases with $k$, so by Lemma~\ref{lem:cluster} its
expectation is at most $1-\E e^{-\lambda\varepsilon N}
=1-\psi(e^{-\lambda\varepsilon})$. This proves \eqref{eq:A1}.

For \eqref{eq:A2}, write $y=\sqrt{1-z}$. Then
$1-\psi(z)=(3y+y^2)/(1+y)^2\le3y$. Take $\varepsilon=\lambda/4$,
which is less than $1/2$ for $\lambda\le1$. Since
$1-e^{-\lambda^2/4}\le\lambda^2/4$, we get
$1-\psi(e^{-\lambda^2/4})\le3\lambda/2$, and the first term of
\eqref{eq:A1} becomes $\lambda[-1/4+\lambda/8+\frac34\log(2/\lambda)]$.
Adding gives \eqref{eq:A2}.
\end{proof}

The same argument applies to an infinite graph of maximum degree at most three.
By Lemma~\ref{lem:cluster} all clusters of $(V,E_{1/2})$ are finite
almost surely. Vertices in different clusters of $E_{1/2}$ do not
interact before time $1/2$, so the process up to time $1/2$ is defined by
running it on each of these finite clusters, and \eqref{eq:A1} and
\eqref{eq:A2} hold for it. This is the process on $\T$ used in
Theorem~\ref{thm:C}.

\begin{proof}[Proof of Corollary~\ref{cor:B}]
On the event $B_r(G_n,o_n)\cong\T_r$, fix an isomorphism and use the same
edge times for corresponding edges of the two balls; all other marks are
independent. Suppose that no vertex of the ball in $G_n$ is frozen at
time $1/2$. Then every edge of the ball opens exactly at its time $U_e$
if $U_e\le1/2$, and every vertex of the ball is warm on $[0,1/2]$. This
is also the history of Aldous's process on $\T_r$ up to time $1/2$, which
freezes nothing then. Applying Theorem~\ref{thm:A} to each of the
$3\cdot2^r-2$ vertices of the ball and taking a union bound gives
\eqref{eq:B1}. If $G_n\to\T$ locally then $b_n(r)\to0$ for each $r$, and
\eqref{eq:B1} tends to zero when $\lambda_n\to0$.
\end{proof}

\section{The freezing-time recursion on the tree}\label{sec:rec}

Remove from $\T$ the edge between the root $o$ and one of its neighbours
$w$. The component of $w$ is a branch. Let
$X$ be the freezing time of $w$ in the process run on this branch
alone, with values beyond $1/2$ replaced by $\infty$. Let $w_1,w_2$ be
the children of $w$, let $X_1,X_2$ be their freezing times in their
own branches, and let $U_1,U_2$ be the times of the edges $ww_1,ww_2$.

\begin{lemma}\label{lem:rec}
Almost surely, $X=\min\bigl(\Lambda_w,\ \{X_i:\ U_i<X_i\}\bigr)$, where
the minimum of an empty set is $\infty$ and values beyond $1/2$ are
replaced by $\infty$ on both sides. The variables $X_1,X_2,U_1,U_2,
\Lambda_w$ are independent and $X_1,X_2$ have the law of $X$. The same
holds for the root $o$ with three children in place of two.
\end{lemma}

Let $U_0$ be the time of the edge $ow$, put $Y=X$ if $U_0<X$ and
$Y=\infty$ otherwise, and define $Y_1,Y_2$ in the same way from
$X_1,U_1$ and $X_2,U_2$. Then Lemma~\ref{lem:rec} reads
$Y=\gamma[U_0]\bigl(\min(\Lambda_w,Y_1,Y_2)\bigr)$, where
$\gamma[u](x)=x$ if $x>u$ and $\gamma[u](x)=\infty$ otherwise. Without
the term $\Lambda_w$ this is the recursion of Aldous's process on the
oriented binary tree, as Ráth, Swart and Terpai write
it~\cite[Section~1.2]{RST21}; the lightning enters only as a third term
of the minimum.

\begin{proof}
First let the branch be a finite tree; all event times are distinct
almost surely. Write $\sigma=\min(\Lambda_w,\{X_i:U_i<X_i\})$.

Up to the freezing time of $w$, each child branch evolves as it
does alone. Indeed, an edge of branch $i$ opens or not according to the
states of its own endpoints, and a strike can change the states in
branch $i$ only by freezing an open cluster meeting branch $i$. Before
$w$ is frozen, such a cluster cannot contain $w$, because $w$ would then
freeze; so it lies in branch $i$, and it is the cluster the strike
freezes in branch $i$ alone. Ordering the event times, this shows by
induction that the two evolutions agree until $w$ freezes.

Now $w$ freezes at its own strike $\Lambda_w$ or when a strike at a
vertex $u$ of some branch $i$ freezes an open cluster containing $w$.
In the second case, at that time $s$ the edge $ww_i$ is open, so $U_i<s$
and $w_i$ was warm at time $U_i$, and $w_i$ belongs to the struck
cluster. Since the evolution of branch $i$ agrees with its own up to
time $s$, $w_i$ freezes at $s$ in branch $i$ alone and was warm before.
Hence $s=X_i>U_i$, and the freezing time of $w$ is at least
$\sigma$.

Conversely, suppose the freezing time of $w$ were larger than
$\sigma$. Then $\sigma\neq\Lambda_w$, so $\sigma=X_i$ for some $i$ with
$U_i<X_i$. At time $U_i$, $w$ is warm, and $w_i$ is warm because its
branch evolves as it does alone and $U_i<X_i$, so the edge $ww_i$ opens.
At time $X_i$, the strike that freezes $w_i$ in branch $i$ alone hits a
vertex $u$ of branch $i$ whose open cluster contains $w_i$. In the joint
process this cluster also contains $w$ through the open edge $ww_i$, so
$w$ freezes at $X_i=\sigma$, a contradiction. If $w$ is already frozen
at time $U_i$, its freezing time is less than $U_i<X_i$, so in
this case too it is at most $\sigma$. Therefore the freezing time
of $w$ equals $\sigma$. The argument does not use the number of children.

For the infinite branch, the clusters of $E_{1/2}$ are finite almost
surely. The values of $X,X_1,X_2$ up to $1/2$, and the event that they
exceed $1/2$, are determined by the marks in the clusters of $w,w_1,w_2$
in $E_{1/2}$. Apply the finite case to a finite subtree containing these
clusters. Independence holds because $X_i$ depends only on the marks in
branch $i$, and the branches are copies of the original one.
\end{proof}

Let $F(t)=\PP(X\le t)$ for $0\le t\le1/2$ and
\[
a(t)=1-\int_{[0,t]}s\,\dd F(s).
\]
Since $X$ is one of the lightning times when it is finite, $F$ has no
atoms. By Lemma~\ref{lem:rec}, $X>t$ if and only if $\Lambda_w>t$ and,
for $i=1,2$, it is not the case that $U_i<X_i\le t$; and
$\PP(U_i<X_i\le t)=\int_{[0,t]}s\,\dd F(s)$. Hence, for $0\le t\le1/2$,
\begin{equation}\label{eq:rde}
1-F(t)=e^{-\lambda t}a(t)^2,\qquad
\PP(\tau_o>t)=e^{-\lambda t}a(t)^3.
\end{equation}

\begin{lemma}\label{lem:ode}
For $\lambda>0$ the function $h(t)=e^{\lambda t}/a(t)-1$ is continuously
differentiable on $[0,1/2]$ and
\begin{equation}\label{eq:ode}
h'=\lambda\,\frac{(1+h)(1+h-t)}{1-2t+h},\qquad h(0)=0.
\end{equation}
Moreover $h(t)\ge e^{\lambda t}-1\ge\lambda t$, and with $H=h(1/2)$,
\begin{equation}\label{eq:exact}
\PP(\tau_o\le\tfrac12)=1-\frac{e^{\lambda}}{(1+H)^3}.
\end{equation}
\end{lemma}

\begin{proof}
Both sides of the first identity in \eqref{eq:rde} are continuous and of
bounded variation. Taking Stieltjes differentials and using
$\dd a=-t\,\dd F$,
\[
\bigl(1-2te^{-\lambda t}a(t)\bigr)\,\dd F(t)=\lambda e^{-\lambda t}a(t)^2\,\dd t.
\]
Since $a\le1$, the factor on the left is at least $1-2te^{-\lambda t}$,
which is positive on $[0,1/2]$ when $\lambda>0$. Hence $F$ has the
continuous density
$F'=\lambda e^{-\lambda t}a^2/(1-2te^{-\lambda t}a)$, and $a'=-tF'$. For
$x=e^{\lambda t}/a=1+h$ this gives
\[
x'=\lambda x+\frac{\lambda t}{1-2t/x}=\lambda x\,\frac{x-t}{x-2t},
\]
which is \eqref{eq:ode}; note $x-2t>0$. Since $(x-t)/(x-2t)\ge1$, we
have $x'\ge\lambda x$ and $x(t)\ge e^{\lambda t}$. Finally
$a(1/2)=e^{\lambda/2}/(1+H)$, and \eqref{eq:exact} follows from the
second identity in \eqref{eq:rde}.
\end{proof}

The derivation uses the independence of the branches, which fails on
graphs with cycles.

\section{Proof of Theorem~\ref{thm:C}}\label{sec:asym}

Constants $C$ below do not depend on $\lambda$, and statements hold for
all sufficiently small $\lambda$. By Theorem~\ref{thm:A} for $\T$ and by
\eqref{eq:exact}, $H\to0$ as $\lambda\to0$; we may assume $H\le1$. Put
\[
D(t)=1-2t+h(t),\qquad I=\int_0^{1/2}\frac{\dd t}{D(t)}.
\]
Writing $1+h-t=\frac12D+\frac12(1+h)$, equation \eqref{eq:ode} becomes
\begin{equation}\label{eq:split}
h'=\frac\lambda2(1+h)+\frac\lambda2\,\frac{(1+h)^2}{D}.
\end{equation}

\emph{A priori bounds.} Since $h\ge\lambda t$ and $h\le H\le1$,
integrating \eqref{eq:split} gives
\[
H\le\lambda+2\lambda\int_0^{1/2}\frac{\dd t}{1-2t+\lambda t}
=\lambda+\frac{2\lambda\log(2/\lambda)}{2-\lambda}
\le C\lambda\log\frac1\lambda .
\]
Since $h\le H$ and $(1+h)^2\ge1$, it also gives
\[
H\ge\frac\lambda2\int_0^{1/2}\frac{\dd t}{1-2t+H}
=\frac\lambda4\log\frac{1+H}{H},
\]
so $H/\lambda\to\infty$.

\emph{The integral $I$.} We show
\begin{equation}\label{eq:I}
I=\frac12\log\frac1H+o(1).
\end{equation}
From \eqref{eq:ode} and $1+h\le2$ we get $h'\le4\lambda/D$, and $D\ge h$
since $t\le1/2$. Let $t_*=1/2-H$ and $h_*=h(t_*)$. On $[t_*,1/2]$ we have
$(h^2)'=2hh'\le8\lambda$, so $H^2-h_*^2\le8\lambda H$ and
$H-h_*\le8\lambda$. As $\lambda=o(H)$ by the a priori bounds above, this gives $h\ge h_*\ge H/2$ on
$[t_*,1/2]$. On $[0,t_*]$ we have $h'\le4\lambda/(1-2t)$, so for
$t\le t_*$
\begin{equation}\label{eq:gap}
0\le H-h(t)\le8\lambda+2\lambda\log\frac{1-2t}{2H}
\le C\lambda\Bigl(1+\log\frac{1-2t}{2H}\Bigr).
\end{equation}
Compare $I$ with
$J=\int_0^{1/2}\dd t/(1-2t+H)=\frac12\log\bigl((1+H)/H\bigr)$:
\[
0\le I-J=\int_0^{1/2}\frac{H-h(t)}{D(t)\,(1-2t+H)}\,\dd t .
\]
On $[0,t_*]$ bound the denominator below by $(1-2t)^2$, use
\eqref{eq:gap} and substitute $u=(1-2t)/(2H)$:
\[
\int_0^{t_*}\frac{H-h(t)}{(1-2t)^2}\,\dd t
\le\frac{C\lambda}{4H}\int_1^\infty\frac{1+\log u}{u^2}\,\dd u
=\frac{C\lambda}{2H}.
\]
On $[t_*,1/2]$, an interval of length $H$, the integrand is at most
$8\lambda/((H/2)H)$, so this part is at most $16\lambda/H$. Hence
$I-J=O(\lambda/H)=o(1)$, and \eqref{eq:I} follows because
$J=\frac12\log(1/H)+O(H)$. In particular $HI\to0$.

\emph{The expansion.} Integrate \eqref{eq:split} over $[0,1/2]$. The
first term contributes $\frac\lambda2(\frac12+O(H))$. In the second,
$(1+h)^2=1+O(H)$, so it contributes $\frac\lambda2(I+O(HI))$. By
\eqref{eq:I},
\begin{equation}\label{eq:H1}
H=\frac\lambda4\Bigl[\log\frac1H+1+o(1)\Bigr].
\end{equation}
Since $\lambda/2\le H$ (take $t=1/2$ in $h\ge\lambda t$, Lemma~\ref{lem:ode})
and $H\le C\lambda L$, we have $\log(1/H)=L+O(\log L)$, and
\eqref{eq:H1} gives $H/(\lambda L)\to1/4$. Therefore
$\log(1/H)=L-\log L+\log4+o(1)$, and by \eqref{eq:H1} again
\begin{equation}\label{eq:H2}
H=\frac\lambda4\Bigl[L-\log L+\log4+1+o(1)\Bigr].
\end{equation}
Finally, $H^2+\lambda H=O(\lambda^2L^2)=o(\lambda)$, so \eqref{eq:exact}
gives
$\PP(\tau_o\le\frac12)=1-(1+\lambda)(1-3H)+o(\lambda)=3H-\lambda+o(\lambda)$.
Inserting \eqref{eq:H2} gives \eqref{eq:C1}, since $1-\frac43=-\frac13$.
\qed

\section{Sharpness and two examples}\label{sec:sharp}

\begin{proof}[Proof of \eqref{eq:sharp}]
The upper bound $\limsup\le3/4$ is \eqref{eq:A2}. For the lower bound,
let $\T_R$ be the ball of radius $R$ in $\T$, rooted at its centre $o$,
and use the same marks on $\T_R$ as on $\T$. The freezing time of
$o$ up to $1/2$ is determined by the marks on the cluster of $o$ in
$E_{1/2}$ (Section~\ref{sec:A}). If this cluster does not reach distance
$R$ from $o$, it is the same in $\T_R$ and in $\T$, and the two
freezing times agree up to $1/2$. Let $q_k$ be the probability that
critical binary branching with $\mathrm{Bin}(2,1/2)$ offspring survives
$k$ generations. Then $q_0=1$ and $q_{k+1}=1-(1-q_k/2)^2=q_k-q_k^2/4$, so
$1/q_{k+1}\ge1/q_k+1/4$ and $q_k\le4/(k+4)$. The cluster of $o$ reaches
distance $R$ with probability at most $3\cdot\frac12 q_{R-1}\le
6/(R+3)$. Hence
\[
\bigl|\PP_{\T_R}(\tau_o\le\tfrac12)-\PP_{\T}(\tau_o\le\tfrac12)\bigr|\le\frac6{R+3},
\]
and letting $R\to\infty$ for fixed $\lambda$ shows that the supremum in
\eqref{eq:sharp} is at least $\PP_\T(\tau_o\le\frac12)$. Theorem~\ref{thm:C}
completes the proof.
\end{proof}

\begin{remark}[The degree bound is needed]
Without a bound on the degree there is no bound of the form of
Theorem~\ref{thm:A}. In a star with $m$ leaves, a leaf whose edge time is
at most $1/4$ and whose lightning time lies in $[1/4,1/2]$ freezes the
centre by time $1/2$, unless the centre is frozen already. These events
are independent over the leaves and each has probability
$p=\frac14(e^{-\lambda/4}-e^{-\lambda/2})$, so the centre is frozen by
time $1/2$ with probability at least $1-(1-p)^m$, which tends to one
when $m=\lceil\lambda^{-2}\rceil$ and $\lambda\to0$.
\end{remark}

\begin{remark}[The first-order term on a fixed graph]\label{rem:cycles}
For a fixed finite graph, two strikes before $1/2$ have probability
$O(\lambda^2)$, and with one strike $v$ freezes exactly when the struck
vertex is in $C_s(v)$ at the strike time $s$; a strike outside $C_s(v)$
freezes a cluster not containing $v$, and $v$ cannot join that cluster
later, because an edge with a frozen endpoint never opens. So
$\frac{\dd}{\dd\lambda}\PP(\tau_v\le\frac12)|_{\lambda=0}
=\int_0^{1/2}\E|C_t(v)|\,\dd t$. For a vertex of a triangle,
$\E|C_t(v)|=1+2t+2t^2-2t^3$ and the integral is $77/96$. The tree has no
such first-order term: there $\E|C_t(o)|=-\frac12+\frac3{2(1-2t)}$ is not
integrable at $1/2$, which is the source of the factor $\log(1/\lambda)$
in \eqref{eq:C1}.
\end{remark}

\section{Numerical values of the constant}\label{sec:num}

We solved \eqref{eq:ode} with an arbitrary-precision Taylor integrator
(\texttt{mpmath.odefun}, 40 digits) for $\lambda=10^{-k}$,
$k=3,4,6,8,10,12$, and evaluated \eqref{eq:exact}. Table~\ref{tab:num}
lists $K(\lambda)=\PP(\tau_o\le\frac12)/(3\lambda/4)-(L-\log L)$, which
by \eqref{eq:C1} tends to $\log4-\frac13=1.05296\ldots$. The convergence
is slow, as expected from the error terms. Write
$\rho(\lambda)=4H/\lambda-\log(1/H)-1$ for the error in \eqref{eq:H1}. Using
\eqref{eq:exact} and the definition of $\rho$ one finds
\[
K(\lambda)-\Bigl(\log4-\frac13\Bigr)
=\rho(\lambda)-\log\frac{4H}{\lambda L}+O(\lambda L^2).
\]
The difference between the third and fourth columns of Table~\ref{tab:num}
is the term $O(\lambda L^2)$; it equals $2.1\cdot10^{-2}$, $3.6\cdot10^{-3}$
and $8\cdot10^{-5}$ for $\lambda=10^{-3},10^{-4},10^{-6}$ and is below
$10^{-5}$ for $\lambda\le10^{-8}$. The
product $L\,\rho(\lambda)$ stays near $1.05$, so $\rho(\lambda)$ decays like
$1/L$, and $\log(4H/(\lambda L))$ is of order $(\log L)/L$; this is the
rate at which $K(\lambda)$ approaches its limit.

\begin{table}[ht]
\centering
\begin{tabular}{rcccc}
\toprule
$\lambda$ & $K(\lambda)$ & $K(\lambda)-\log4+\frac13$ & $\rho(\lambda)-\log\frac{4H}{\lambda L}$ & $L\,\rho(\lambda)$\\
\midrule
$10^{-3}$ & $1.11663$ & $0.06367$ & $0.08455$ & $1.1023$\\
$10^{-4}$ & $1.13639$ & $0.08343$ & $0.08702$ & $1.0511$\\
$10^{-6}$ & $1.14003$ & $0.08706$ & $0.08714$ & $1.0508$\\
$10^{-8}$ & $1.13468$ & $0.08172$ & $0.08172$ & $1.0544$\\
$10^{-10}$ & $1.12848$ & $0.07552$ & $0.07552$ & $1.0539$\\
$10^{-12}$ & $1.12276$ & $0.06980$ & $0.06980$ & $1.0520$\\
\bottomrule
\end{tabular}
\caption{The tree constant from the differential equation \eqref{eq:ode}.}
\label{tab:num}
\end{table}

\section*{Acknowledgments and funding}

The research presented in this paper was supported by the European
Research Council (ERC) under the European Union's Horizon Europe research
and innovation programme (grant agreement No.~101041711) and by the Simons
Foundation as part of the Collaboration on the Mathematical and Scientific
Foundations of Deep Learning. It was also supported by Heights Labs, by
Convex Nexus Capital and by the Israel Science Foundation (ISF grants
2258/19 and 4101/25).

This work was done by multiple AI agents, Claude (Anthropic) and Codex
(OpenAI), connected to Math Brain~\cite{MathBrainRepo} and working
together inside the Brain framework~\cite{BrainFramework}. The method is
described in~\cite{MathBrainPaper}.

\section*{Declaration of competing interest}

The author holds equity in Heights Labs and in Convex Nexus Capital and
has a management role at Convex Nexus Capital.

\section*{Declaration of generative AI and AI-assisted technologies}

This paper is part of a broader research project on proving open
problems with AI systems, carried out with Claude (Anthropic) and Codex (OpenAI) connected
to Math Brain~\cite{MathBrainRepo} inside the Brain
framework~\cite{BrainFramework}, as described in~\cite{MathBrainPaper}.
Codex (model GPT-6.1-sol) proposed and proved the results of
Sections~\ref{sec:A}--\ref{sec:sharp}, and a separate Codex run audited
them. Claude (Opus~5.5) checked the proofs, wrote the computation of
Section~\ref{sec:num}, and wrote the text. Frozen percolation on the
three-regular tree is due to Aldous~\cite{Ald00}; the finite-graph model
with lightning and Problem~2 are due to Ráth, Swart and
Terpai~\cite{RST21}. Without lightning, the recursion of
Lemma~\ref{lem:rec} is the recursion of Aldous's process as Ráth, Swart
and Terpai write it~\cite{RST21}. Recursive distributional equations of the kind
used in Section~\ref{sec:rec} are surveyed by Aldous and
Bandyopadhyay~\cite{AB05}, and the complete-graph model with lightning
and its window curve, which Theorem~\ref{thm:C} parallels, are due to
Ráth, Szőke and Yeo~\cite{RSY26}. These results are theirs. The author directed the work and is responsible for
the content of the paper and for decisions concerning its dissemination.

\section*{Data and code availability}

The script \texttt{residual\_check.py} and its output are available at
\url{https://github.com/yspennstate/frozen-percolation-rare-lightning}.

\begin{thebibliography}{99}

\bibitem{Ald00}
D.~J. Aldous, \emph{The percolation process on a tree where infinite
clusters are frozen}, Math. Proc. Cambridge Philos. Soc. \textbf{128}
(2000), 465--477, doi:10.1017/S0305004199004326.

\bibitem{AB05}
D.~J. Aldous and A.~Bandyopadhyay, \emph{A survey of max-type recursive
distributional equations}, Ann. Appl. Probab. \textbf{15} (2005),
1047--1110, doi:10.1214/105051605000000142.

\bibitem{BKN12}
J.~van den Berg, D.~Kiss and P.~Nolin, \emph{A percolation process on the
binary tree where large finite clusters are frozen}, Electron. Commun.
Probab. \textbf{17} (2012), no.~2, 1--11, doi:10.1214/ECP.v17-1694.

\bibitem{Rat09}
B.~Ráth, \emph{Mean field frozen percolation}, J. Stat. Phys.
\textbf{137} (2009), 459--499, doi:10.1007/s10955-009-9863-5.

\bibitem{RST21}
B.~Ráth, J.~M. Swart and T.~Terpai, \emph{Frozen percolation on the
binary tree is nonendogenous}, Ann. Probab. \textbf{49} (2021),
2272--2316, doi:10.1214/21-AOP1507.

\bibitem{RSY26}
B.~Ráth, M.~Szőke and D.~Yeo, \emph{The window process of slightly
subcritical frozen percolation}, preprint, 2026, arXiv:2610.37320.

\bibitem{RT09}
B.~Ráth and B.~Tóth, \emph{Erd\H{o}s--Rényi random graphs + forest fires =
self-organized criticality}, Electron. J. Probab. \textbf{14} (2009),
1290--1327, doi:10.1214/EJP.v14-653.

\bibitem{MathBrainRepo} Y.~Shmalo, \emph{Math Brain}, public repository, 2026, \url{https://github.com/yspennstate/math-brain-public/tree/8c9a6bd478c4ec4082347151b7ae16db50c917dd}.

\bibitem{MathBrainPaper} Y.~Shmalo, \emph{Math Brain: shared, auditable memory for multi-agent mathematical research}, October 2026, doi:10.5281/zenodo.24150837,
\url{https://github.com/yspennstate/math-brain-public/tree/main/papers/math-brain-2026-10-24}.

\bibitem{BrainFramework} Y.~Shmalo, \emph{Brain Framework}, public software repository, 2026, \url{https://github.com/yspennstate/brain-framework-public/tree/5629e644cdd510173ae97e9aa92642dabd22dc57}.

\end{thebibliography}

\end{document}
